\section{Possible Final State}
\label{sec:final-state}

This section describes what the project could demonstrate at the end of the
course. It is a target to steer by rather than a commitment. The work starts
small and grows in stages (Section~\ref{sec:stages}) towards the target system
in Section~\ref{sec:target}.

\textbf{Answers to the problem statement.} Detection uses the drones' own
cameras and a two-tier classifier (Section~\ref{sec:target-approach}); the
detection ranges and the number of drones needed follow from
Section~\ref{sec:derived}. Interception compares a planning-based controller
with one based on reinforcement learning on the same scenarios
(Section~\ref{sec:target-approach}).

\subsection{Stages}
\label{sec:stages}

Table~\ref{tab:stages} lists the stages from the current state to the final
one. Later stages only add to earlier ones, so the system works end to end
throughout the project.

\begin{table}[H]
\centering
\small
\begin{tabular}{|p{0.22\textwidth}|p{0.70\textwidth}|}
\hline
\textbf{Stage} & \textbf{What works} \\
\hline
0. Current state &
One interceptor follows one hostile drone at a fixed altitude of 50\,m,
given a noisy estimate of its position, and keeps a safe distance behind it. \\
\hline
1. Lock-on &
A narrower gimbal camera than the platform's (Section~\ref{sec:camera-choice}); a simple detector for a dark object against the sky;
a tracker on the camera bearing; a locked state. One hostile drone against
the sky, starting inside the camera's view. \\
\hline
2. Search and intercept &
The narrow camera scans the outward sector; hostile drone at $v_h$ and
interceptor at $v_i$; meeting point on the hostile drone's path and
proportional navigation once locked; a ring of $N$ drones, with the base
station assigning the interceptor. \\
\hline
3. Classification &
Image crops sent to the base station; a classifier over image sequences;
birds as distractors; an \emph{investigate} role; tuning so that no drone is
missed. \\
\hline
4. Wide coverage &
A second, wide camera on each drone's gimbal and slew-to-cue; a convolutional neural network
detector that also works against trees and ground. \\
\hline
5. Learned controller (stretch goal) &
An interception policy trained with reinforcement learning, compared with the
planning-based controller; several hostile drones and decoys. \\
\hline
\end{tabular}
\caption{Stages from the current state to the final state.}
\label{tab:stages}
\end{table}

Every stage is evaluated as described in Section~\ref{sec:evaluation}. In
addition, the derived values in Section~\ref{sec:derived} are checked against
the simulation, for example whether the intercept margin in
Table~\ref{tab:feasibility} matches the measured one.

\subsection{Limitations}

Everything runs in simulation; no real drones or sensors are used. The
derived values rest on simplifications: a straight, radial approach at
constant speed, detection as soon as the target is $p_{\min}$ pixels wide, a
clear sky behind the target, and one hostile drone at a time. Real detection
is probabilistic and weaker against clouds, trees or ground. What an intercept
ends with, for example escorting or neutralising the hostile drone, is left
open: the final state ends when the interceptor is locked on and within
intercept distance.
